\chapter{An independent uniform NS--PC comparison}\label{tx:sec:comparison}
The linear PC lower bounds were proved by transporting the PC closure
itself. We now establish a different result: every bounded-degree PC
consequence has an NS identity of degree bounded independently of
universe size. The comparison applies to arbitrary polynomials, not
only invariant targets or refutations. Its qualitative consequence is
that NS and PC assign every fixed uniform request the same type.
Throughout, $\Gamma_n$ is a fixed uniform template over a field $K$
of fixed prime characteristic $p$. Recall that $I_d(n)$ is the span of
axiom products of total-product degree at most $d$, while $C_d(n)$ is
the ordinary-PC consequence closure inside $P_{n,\le d}$.

\section{What comparison can and cannot prove}
Suppose you already know that NS degree is large. It is tempting to use
$h_{\PC}\le h_{\NS}$ as a PC lower bound. The inequality points in the
wrong direction. Even an upper comparison $h_{\NS}\le F(h_{\PC})$
need not preserve a linear scale if $F$ is large. Our linear PC bound
therefore came from its own closure transport, not from inverting $F$.

The purpose of comparison is instead to say that bounded PC degree can
never coexist with unbounded NS degree on the same sequence of sizes.
This is its role in the running example of
Chapter~\ref{tx:ch:running-order}: the elementary induction-chain
restriction supplies a growing NS obstruction, and comparison rules
out recurring bounded PC degree for the original uniform order
family. The linear PC constant then comes from PC's own lifting
range, without inverting the possibly large function $F$.
It also controls recurring finite levels well enough to identify the
same one of the four types. The proof will make this statement for
\emph{every} polynomial vector in the PC consequence space, before
specialising to any uniform target.

\section{Why one plateau makes a space closed}
The ideal generated by $\Gamma_n$ consists of finite sums $\sum_f a_ff$
with arbitrary polynomial multipliers. The spaces $I_e(n)$ impose a
product-degree bound on those sums, so they need not be ideals:
multiplication may leave the allowed degree window. The argument below
finds enough closure inside a fixed polynomial ambient without asserting
that a truncated NS space is closed under arbitrary multiplication.

Imagine a chain of NS consequence spaces with increasing allowed degree,
but look only at their intersection with the fixed degree-$d$ ambient.
Thus $I_{d+j}(n)\cap P_{n,\le d}$ allows larger certificates while
retaining only consequences of polynomial degree at most $d$.
An invariant-chain bound forces two successive intersections to be equal.
At that equality, multiplication of a degree-$d-1$ vector cannot create
a new degree-$d$ vector: it has only used one extra unit of NS certificate
degree, precisely the step covered by the plateau. Thus the plateau is
closed under the ordinary-PC rules.

This is the core idea. The first lemma supplies the uniform chain bound;
the second argument turns a plateau into whole-consequence containment.
Neither step assumes symmetry of the proof lines or the target.

The length being bounded is the number of strict inclusions in an
ascending chain. For example, $A_0\subsetneq A_1=A_2\subsetneq A_3$ has two strict
inclusions. Invariance means that relabelling sends each subspace to
itself. Exact deletion transports the inclusions, and the
detector keeps a strict inclusion strict by preserving its nonzero
quotient. This is why a dimension bound at small rank can control
chains at every rank.

\begin{lemma}[Uniform chain bound]\label{tx:lem:chains}
For any bounded basis system $\mathcal U$, strict $S_n$-invariant
subspace chains in $K[\mathcal U_n]$ have a length bound independent of $n$.
\end{lemma}
\begin{proof}
Let $r$ bound the basis support. Over an algebraic closure, choose a
fixed $q=p^a>2r$. At $N=n+q$ with $n\ge4r+1$, deletion is exact at the
doubled budget $2r$. Lemma~\ref{tx:lem:detect} shows that each nonzero
subquotient $B/A$ of $K[\mathcal U_N]$ survives. Thus
$A\subsetneq B$ implies
$\mathcal D(A)\subsetneq\mathcal D(B)$: the whole strict chain becomes
a strict chain of the same length in $K[\mathcal U_n]$.

Repeatedly subtract $q$ while the rank exceeds $M=4r+q$. Every retained
rank in such a step exceeds $4r$, so all steps are eligible. The final
chain has at most $\dim K[\mathcal U_M]$ strict inclusions, because a
smaller-rank basis injects into the rank-$M$ basis. The same bound covers
the initial ranks at most $M$. Each step may choose its own generic
coefficients: the detector applies to arbitrary invariant subspaces,
without a naturality assumption on the chain across ranks. Finally,
scalar extension preserves strict inclusions, giving the result over
the original coefficient field.
\end{proof}

Only a finite exactness threshold at each fixed support is needed for
this argument. In characteristic two the earlier quadratic construction also suffices.
For an arbitrary positive characteristic the completed linear construction
already provides the needed threshold. The
bound may be large; its role is uniformity in $n$, not a sharp comparison
of the two minimum degrees.

\begin{theorem}[Uniform comparison for all consequences]\label{tx:thm:comparison}
For every uniform axiom template over a field of prime characteristic
$p$, there is a nondecreasing $F:\NN\to\NN$, independent of $n$ and of
the target, with $F(d)\ge d$ and
\begin{equation}\label{tx:eq:whole-comparison}
                              C_d(n)\subseteq I_{F(d)}(n)
                       \qquad(n,d\ge0).
\end{equation}
Consequently, every polynomial with an ordinary-PC derivation of degree
at most $d$ from $\Gamma_n$ has a total-product NS identity of degree at
most $F(d)$ from the same axioms. For the aggregate profiles of any
stable uniform target family, and for the profiles of any one fixed
target orbit,
\begin{equation}\label{tx:eq:target-comparison}
                      h_{\PC}(n)\le h_{\NS}(n)\le F(h_{\PC}(n)),
                      \qquad F(\infty)=\infty.
\end{equation}
Their NS and PC types in Theorem~\ref{tx:thm:classification} therefore agree.
\end{theorem}
\begin{proof}
Fix $d$. Monomials of degree at most $d$ form a permutation basis of
support at most $kd$. Let $L_d$ bound strict invariant-subspace chains in
$P_{n,\le d}$, uniformly in $n$, by Lemma~\ref{tx:lem:chains}.
Consider the invariant spaces
\[
              W_j=I_{d+j}(n)\cap P_{n,\le d},
                    \qquad 0\le j\le L_d+1.
\]
They form an ascending chain. Some $W_j=W_{j+1}$ with $j\le L_d$,
otherwise there would be more than $L_d$ strict inclusions.

This plateau contains every axiom usable within degree $d$ and is closed
under linear combinations. If $f\in W_j$ and $\deg f\le d-1$, then
multiplication of an NS identity for $f$ by a variable gives
$x_vf\in I_{d+j+1}(n)$. Its degree is at most $d$, hence
$x_vf\in W_{j+1}=W_j$. For $d=0$ there is no permitted multiplication
of a nonzero line; the zero line already lies in $W_j$.
Induction on \emph{any} degree-$d$ PC derivation now places every one of
its lines in $W_j$. In particular,
\[
                          C_d(n)\subseteq I_{d+L_d}(n).
\]
This induction retains cancellations: a lower-degree linear combination
can be multiplied because it lies in the plateau, even if its separate
summands have degree $d$.

Choose
\[
                          F(d)=d+\sum_{u=0}^d L_u.
\]
It is nondecreasing, satisfies $F(d)\ge d$, and gives
\eqref{tx:eq:whole-comparison}. No symmetry of a target was used.
Conversely, expansion of NS identities into monomial--axiom products
gives a same-degree PC derivation, so the first inequality in
\eqref{tx:eq:target-comparison} holds. If the aggregate PC degree is finite,
every requested target belongs to $C_d(n)$ at that budget; the inclusion
puts every target in $I_{F(d)}(n)$. Empty requests have degree zero.
Infinite PC degree forces infinite NS degree by the first inequality.
These observations prove the complete extended-order comparison.

For type agreement write $g=h_{\PC}$ and $h=h_{\NS}$.
The inequalities $g\le h\le F(g)$ transfer eventual boundedness in both
directions. If $g$ takes one finite value $d$ infinitely often, then on
those sizes $h$ lies in the finite interval $[d,F(d)]$; some finite value
of $h$ recurs. A recurring finite value of $h$ likewise forces a recurring
finite value of $g$. Thus the recurring sets are empty together.
If recurring values of $g$ are unbounded, the same finite-interval
argument yields unbounded recurring values of $h$. Conversely, choose a
recurring value of $h$ exceeding $F(B)$. On its infinite level set,
$g>B$ and $g\le h$, so some finite value above $B$ recurs for $g$.
Thus infinitude of the recurring sets agrees as well. Boundedness,
emptiness and infinitude determine exactly the four alternatives.
\end{proof}

\begin{corollary}[Common finite peaks]\label{tx:cor:common-peaks}
If a fixed uniform request has common type T4, there is a nondecreasing
$F$ such that, beyond every bound $B$, some size $n\ge B$ has
\[
 h_{\PC}(n)=d\ge B,\qquad d\in E(h_{\PC}),\qquad
 d\le h_{\NS}(n)\le F(d)<\infty,
\]
with $n\le(p+2)(\kappa d+\sigma)$. Both degrees satisfy the rational
linear lower bound in \eqref{tx:eq:finite-peaks} at these same sizes.
Here $E(h)$ is the set of recurring finite degrees, and
$R(d)=\kappa d+\sigma$ is the applicable transfer budget: $(\kappa,\sigma)
=(2k,s+t)$ for indexed targets and $(k,s)$ for an invariant target.
\end{corollary}
\begin{proof}
Take the finite PC peaks from Theorem~\ref{tx:thm:classification}.
The pointwise lower comparison gives the NS lower bound, and the uniform
upper comparison gives its finiteness at those sizes.
\end{proof}

The comparison does not identify exact degrees, easy-class values or
recurring degree sets. Its potentially large function $F$ cannot be
inverted to obtain a linear PC lower bound from a linear NS lower bound.
That quantitative result came directly from the PC transfer theorem.

\section{Reading the result as three tests of a profile}
A profile has a recurring finite level if it takes one finite value at
infinitely many sizes. The comparison sends that infinite set into a
finite interval of degrees in the other system, so some value of the
other profile repeats infinitely often. It also preserves whether
recurring levels are unbounded: choose levels above $F(B)$ and use
pigeonhole inside a finite interval. Finally, it preserves eventual
boundedness directly. These three tests distinguish T1, T2, T3 and T4.

This is why one does not need equality of exact degrees, equality of
recurring levels, or equal easy-class values to obtain same-type
agreement. For common finite T4 peaks, take a finite PC peak; pointwise
PC\,$\le$\,NS supplies its lower bound, and the upper comparison supplies
finiteness of the NS degree at that same size.
